% Short running form: see the note on head width in format.tex.
\runninghead{Seek the Lord While He May Be Found}
\properlane{conversion}{Seek the Lord While He May Be Found: The Day as a Human Life}
\label{sec:seek-the-lord}

The same day can be measured on a smaller clock. It is one human life, from
the morning of childhood to the last hour of old age, and the householder goes
out into every part of it. God calls at every age; the one who comes late is
truly received; therefore no one may despair. And because no one knows how
long his day will be, no one may wait for a later hour. Heard in this way the
Mass is a summons to turn now, addressed equally to those who have served long
and to those who have long stood idle. The prophet's ``while he may be found''
is its motto, and nearly every other text supplies either the promise that a
late cry is heard or the warning that the time is short.

\subsection{The hours of a life}

Chrysostom states this sense as the plain meaning of the figures. ``By a
vineyard He meaneth the injunctions of God and His commandments: by the time
of laboring, the present life: by laborers, them that in different ways are
called to the fulfillment of the injunctions: by early in the morning, and
about the third and ninth and eleventh hours, them who at different ages have
drawn near to God, and approved themselves'' (\work{Homilies on Matthew}
64.3). Jerome gives it as his own view, \latin{Mihi videntur}, and fills the
hours with names and ages. The labourers of the first hour are Samuel,
Jeremiah and John the Baptist, who can say with the psalmist, ``From my
mother's womb thou art my God'' (Ps 22 (21):11); those of the third began to serve God from
puberty; those of the sixth took up Christ's yoke in mature age; those of the
ninth when already declining toward old age; those of the eleventh in extreme
old age, \latin{et tamen omnes pariter accipiunt praemium, licet diversus
labor sit}: and yet all alike receive the reward, though the labour differs
(\work{Commentarii in Matthaeum} III, PL 26, 141).

Augustine and Gregory, who also read the hours as the ages of the world, add
this sense as a second one and say so. ``Besides that solution of the
parable,'' Augustine writes, ``the parable may be seen to have an explanation
in respect even of this present life'': infants fresh from the womb are
called at the first hour, boys at the third, young men at the sixth, those
verging on old age at the ninth, the altogether decrepit at the eleventh,
``yet all these are to receive the one and the same denarius of eternal life''
(\work{Sermo} 87.7). Gregory follows the sun through the day. Childhood is the
morning of our understanding; adolescence is the third hour, when the sun
climbs and the heat of age grows; youth is the sixth, when the sun stands at
the centre and strength is at the full; old age is the ninth, when the sun
comes down from the height; and the eleventh is the age called decrepit
(\work{Homiliae in Evangelia} 19.2). Aquinas sets the same sense beside the
ages of the world: the life of a man is one day, childhood is its morning, and
some are given no occasion of returning to God until old age (\work{Super
Matthaeum} c.~20, lect.~1).

Gregory also says what it is to work in this vineyard. Those labour for the
Lord who think \latin{non sua, sed lucra dominica}, not of their own gains but
of the Lord's, who serve with zeal of charity, watch to win souls and hasten
to bring others with them to life. \latin{Nam qui sibi vivit}, for he who
lives for himself and feeds on the pleasures of his flesh is rightly rebuked
as idle, because he pursues no fruit of God's work (19.2). Idleness in the
market-place is thus no absence of activity. It is a life whose activity has
no reference to God.

\subsection{What the parable is for}

Chrysostom asks why the Lord told the parable when he did. He had just spoken
to the rich young man and to Peter about the casting away of riches, which
``needed much vigor of mind and youthful ardor''; so, ``in order to kindle in
them a fire of love, and to give vigor to their will, He shows that it is
possible even for men coming later to receive the hire of the whole day''. The
parable has two audiences and one lesson for each: it is spoken ``to the
former, that they may not be proud, neither reproach those called at the
eleventh hour; to the latter, that they may learn that it is possible even in
a short time to recover all'' (64.4). Its object is ``to render more earnest
them that are converted and become better men in extreme old age, and not to
allow them to suppose they have a less portion'' (64.3).

On this interpretation the murmuring cannot be a fact about the saved, and Chrysostom
says so with force: ``There is no one who is thus justifying himself, or
blaming others in the kingdom of Heaven; away with the thought! for that place
is pure from envy and jealousy.'' The complaint is a device of the story. The
Lord brings in others displeased at the late comers' blessings ``to teach us
that these have enjoyed such honor, as could even have begotten envy in
others'', as a man might say, ``Such a one blamed me, because I counted thee
worthy of much honor'', when no one had blamed him at all. Why, then, were not
all hired at dawn? ``As far as concerned Him, He did hire all; but if all did
not hearken at once, the difference was made by the disposition of them that
were called.'' God called Paul when Paul was ready to obey, and so too the
thief, ``although He was able to have called him even before, but he would
not have obeyed'' (64.3). That sentence returns in the third interpretation.

Gregory finds the same thief in the parable. He came at the eleventh hour,
and \latin{etsi non habuit per aetatem, habuit tamen sero per poenam}: his
lateness was not of age but of punishment. He confessed God on the cross and
breathed out the spirit of life almost with the words. And the householder began to pay from
the last, because he brought the thief to the rest of paradise before he
brought Peter (19.3).

\subsection{Certain of the reward, uncertain of the day}

If the last hour earns the whole wage, why not wait for it? Augustine raises
the objection himself: ``lest any put off to come into the vineyard, because
he is sure, that, come when he will, he shall receive this denarius.'' His
first answer comes from the story. None of the labourers bargained for a later
start. Those found at the third hour did not say, ``Wait, we are not going
thither till the sixth hour'', for what the householder would give ``was in
the secret of His own counsel: do thou come when thou art called''. His second
answer comes from life. To the man in his vigour who proposes to come at the
eleventh hour he says: ``The Householder hath it is true promised thee a
denarius, if thou come at the eleventh hour, but whether thou shalt live even
to the seventh, no one hath promised thee \ldots\ Why then dost thou put off
him that calleth thee, certain as thou art of the reward, but uncertain of the
day? Take heed then lest peradventure what he is to give thee by promise, thou
take from thyself by delay'' (87.8).

He then names the two opposite ways of missing the call, for ``these two
things are the death of souls, despair, and perverse hope''. Some reckon their sins
too great to be forgiven and give themselves up to them; to these the
householder goes out through Ezekiel, ``In whatsoever day a man shall turn
from his most wicked way, I will forget all his iniquities.'' Others reason
from that very promise, ``why should I turn to-day and not to-morrow?'', and to
these he goes out again by another book: ``Make no tarrying to turn to the
Lord, and put not off from day to day.'' God has promised forgiveness to the
one who turns, ``but no one has promised thee to-morrow'' (87.10--11). Gregory
presses the same point on the old: \latin{saltem in ultima aetate
resipiscite}, at least in your last age come to your senses, and come late to
the ways of life now that you have little labour left to give (19.3). From the
verse's closing sentence on the many called and the few chosen, which this
Mass does not read, he draws two further cautions: no one should presume on
himself, and no one should despair of a neighbour sunk in vice, \latin{quia
divinae misericordiae divitias ignorat}, because he does not know the riches
of the divine mercy (19.6).

\subsection{While he may be found}

The first reading says in the prophet's words what Augustine says in the
preacher's. Jerome's comment on ``Seek ye the Lord, while he may be found''
marks the limit of the time: seek him \latin{dum estis in corpore, dum datur
locus poenitentiae}, while you are in the body, while room is given for
repentance; and seek him \latin{non loco, sed fide}, not by place but by
faith. Seeking is not enough, he continues, unless the impious leaves his
former ways; but to the one who does return God is not slow. He draws near to
those who draw near to him, \latin{et filio longo post tempore revertenti,
laetus occurrit}: he runs gladly to meet the son who comes back after a long
time (\work{In Isaiam} XV, PL 24, 553). Aquinas divides the passage into four
steps, and they are the steps of a conversion: the fit time, \latin{antequam
adversitas vel mors veniat}, before adversity or death comes; the manner, in
which the \latin{impius} turns from sin against God and the \latin{iniquus}
from sin against his neighbour; the fruit, which is manifold mercy; and the
removal of mercy's obstacle, since God's thoughts are of mercy where ours are
of requital (\work{Super Isaiam} c.~55). His two directions of sin answer to the two
loves on which the Collect rests the law. To forsake one's way is to come back
to both.

The psalm, in Augustine's hands, turns into the same sermon. At ``patient and
plenteous in mercy'' he asks how many sinners God puts up with, \latin{non
damnans, sed exspectans}, not condemning but waiting, and then he turns on his
hearer: \latin{vocat te nunc, exhortatur te nunc; exspectat donec resipiscas,
et tu tardas!} He calls you now, he exhorts you now; he waits until you come
to your senses, and you delay! It is part of God's great mercy, he adds, that
he has made the day of our death uncertain, so that, expecting daily to
depart, we may at some time be converted (\work{Enarrationes in Psalmos}
144.11, PL 37, 1876). The argument is the one he makes from the parable in
\work{Sermo} 87, here drawn from the psalm's own verse. Bellarmine glosses the
same word: God is patient, ``which means that his mercy is continuous; for no
matter how often we may provoke him, he will not turn to anger at once, but
rather waits to see would we do penance''. The strophes on either side supply
the two halves of this interpretation. ``Every day will I bless thee'' is the
voice of the labourer who does not wait for a better day, and ``The Lord is
nigh unto all them that call upon him'' is the promise made to the one who
calls at the eleventh hour.

\subsection{A life that bears fruit}

Paul, in the second reading, is a labourer near the evening of his day, and he
does not count his hours. Chrysostom notices how exactly he weighs the present
life. By saying ``to die is gain'' he casts out the desire of it; by saying
that to live in the flesh is ``the fruit of my work'' he shows ``that the
present life also is needful, if we use it as need is, if we bear fruit; since
if it be unfruitful, it is no longer life. For we despise those trees which
bear no fruit''. And he adds the sentence that joins the letter to the
parable's market-place: ``God hath made thee live, that thou mayest live to
Him'' (\work{Homilies on Philippians}~3). Paul is like the wrestler who might
be crowned and chooses to go on contending, and his reason is the vineyard's:
``that I may make the plot which I have planted bear much fruit'' (\work{Homilies
on Philippians}~4). Aquinas reads verse~22 in the same sense: if life brings
the fruit that Christ is magnified, then life in the flesh \latin{est bona, et
fructuosa}, is good and fruitful (\work{Super Philippenses} c.~1, lect.~3).
Gregory's definition of the labourer, one who seeks the Lord's gains and not
his own, is a description of the Paul of this passage. The comparison between
Paul and the labourers is the editor's and none of these authors', but each of
them measures a life by its fruit for God.

The second reading's last line, ``Only let your conversation be worthy of the gospel
of Christ'', is the point at which Chrysostom's two homilies meet. His sermon
on the parable ends by demanding ``a life suitable to our faith'', since all
the Lord's parables ``demand virtue in our works'' (\work{Homilies on Matthew}
64.4); his sermon on the letter hears in Paul's ``only'' that this and nothing
else is to be sought. The late comer is received for nothing, and what is then
asked of him is a whole life.

\subsection{From the opened heart to the altar}

Conversion begins in hearing, and the Alleluia verse asks for the hearing. In
the opening lines of his sermon on this parable Augustine describes how God
cultivates us as a husbandman his field: he does not cease ``to open our
heart, as it were, by the plough of His Word, to plant the seed of His
precepts, to wait for the fruit of piety'' (87.1). Lydia is the figure of that
cultivation. Her heart is opened, she attends to what is said, she is baptized
with her household, and at once she opens her house (Acts 16:14--15). The
assembly that asks for her grace before the Gospel asks to be made able to
answer the householder when he comes out.

The Missal's texts give the turning its beginning and its end. At the
entrance the Lord promises to hear the people from whatever tribulation they
cry out of, and that promise is the ground on which a late cry, even the
thief's, is not refused. The Collect asks that those who keep God's precepts
may come to eternal life; since Chrysostom's vineyard is ``the injunctions of
God and His commandments'', the prayer asks in plain words for what the
parable pictures, labour in the vineyard and the denarius at evening. The
Prayer over the Offerings asks that what is professed by faith be laid hold of
in the sacraments, and the turning that began in hearing is completed at the
altar. The first Communion antiphon is the convert's own wish, ``O! that my
ways may be directed to keep thy justifications'', with Hilary's warning
against doing anything \latin{dissoluto animo}, with a slack mind. The second
is Gregory's examination of conscience. Hearing ``I know mine, and mine know
me'', he tells his people, \latin{Videte si oves eius estis}: see whether you
are his sheep, see whether you know him, and know him \latin{non per fidem,
sed per amorem \ldots\ non ex credulitate, sed ex operatione}, not by faith
but by love, not by belief but by deed; for ``He who saith
that he knoweth him and keepeth not his commandments is a liar'' (1~Jn 2:4;
\work{Homiliae in Evangelia} 14.4). The Prayer after Communion asks for that
same proof, that the effect of redemption be possessed \latin{et moribus}, in
conduct as well as in the mysteries.

\subsection{The reward of delay, and the silence about the murmurers}

Augustine names the chief difficulty himself: the parable, taken alone, seems
to reward delay. His answer, that the wage is certain and the day is not, is
true and comes from outside the story, since within it every labourer lives
until evening. The second difficulty is one of proportion. The murmuring and
the householder's reply occupy a third of the parable, and on this
interpretation they have little to do; Chrysostom says outright that they are
not to be pressed. An account of the parable that leaves its ending idle
invites another that begins there. The Entrance Antiphon, the Collect, the
Prayer over the Offerings and the Prayer after Communion take their places
here by their own wording; the placing is the editor's, not a
commentator's.

\subsection{The four senses of this interpretation}

\begin{fourSenses}
\item[Literal.] Men are hired at five hours of one working day and all are
paid the day's wage. A prophet calls his people to seek the Lord and return to
him while there is time, promising pardon. An apostle who longs to be with
Christ judges that his remaining alive is needful for others, and asks of them
only a life worthy of the Gospel.

\item[Allegorical.] The hours of the day are the ages of a life, from the
child called from the womb to the old man at the edge of night (Jerome;
Augustine; Gregory). The householder who goes out is Christ making himself
known (Augustine alone, \work{Sermo} 87.9). The thief who confessed God on the
cross is the labourer of the eleventh hour, paid before the first (Gregory).
The God of
the first reading is the father who runs gladly to meet the son returning
after a long time (Jerome).

\item[Moral.] Turn today. Do not despair over a wasted past, for all is
forgiven the one who turns; do not presume on a future no one has promised
(Augustine). Work for the Lord's gains and not your own (Gregory), let life
bear fruit (Chrysostom; Aquinas), prove by deeds that you know the shepherd
(Gregory), and despair of no one else.

\item[Anagogical.] The evening is each person's death, whose day God has
mercifully hidden (Augustine), and beyond it the common reckoning when the
steward calls the labourers. What was sought in time is then possessed: to be
dissolved and to be with Christ, which Paul calls ``by far the better'', and
the one denarius of eternal life.
\end{fourSenses}
